% Section 2: Literature Review. Shared by both templates.
\section{Literature Review}\label{sec:lit}

\subsection{Foundations}
Gode and Sunder showed that agents with no intelligence beyond a budget constraint already produce
near-efficient allocations in a double auction, so much of market behaviour comes from the trading
mechanism itself~\cite{gode1993}. Chiarella and Iori added fundamentalist, chartist and noise
traders interacting through an explicit order book~\cite{chiarella2002}; Preis et al.\ showed that a
simple multi-agent order book with price--time priority reproduces fat tails once order placement
depth depends on the trend~\cite{preis2006}; and Farmer et al.\ found that zero-intelligence order
flow explains much of the observed spread and price impact~\cite{farmer2005}. Cont's stylised facts
(heavy tails, volatility clustering and no linear autocorrelation) are the standard checklist for
validating such models~\cite{cont2001}. On the computing side, Lysenko and D'Souza built the first
largely GPU-resident ABM framework~\cite{lysenko2008}, Fujimoto surveyed parallel discrete-event
simulation~\cite{fujimoto1990}, and Amdahl~\cite{amdahl1967}, Gustafson~\cite{gustafson1988},
Karp and Flatt~\cite{karp1990} and Williams et al.~\cite{williams2009} give the performance models
used in this report.

\subsection{Recent agent-based market simulators}
\textbf{ABIDES-Gym}~\cite{amrouni2021} (ACM ICAIF 2021) wraps ABIDES~\cite{byrd2020abides}, a
message-driven discrete-event simulator with a NASDAQ-like exchange and price--time priority, as an
OpenAI Gym environment for reinforcement learning. It offers high fidelity and realistic latency,
but its kernel is a single-threaded event loop and the paper reports no run times; a later study
cites ABIDES at about 29{,}000 messages per second~\cite{mascioli2024}.

\textbf{JAX-LOB}~\cite{frey2023} (ACM ICAIF 2023) is the first GPU-accelerated LOB simulator. Each
book is a fixed-capacity array, and matching runs inside a JAX loop. The authors note that
processing messages within one book is inherently serial, so all parallelism comes from running
thousands of independent books with \texttt{vmap}: with 1000 books, a message costs about
2.6\,$\mu$s per book on an RTX 2080 Ti, and RL training is about 7$\times$ faster. Its limits are
pre-allocated capacity, a cost set by the slowest branch, and replayed order flow that does not
react to the agent. It does not measure the end-to-end serial fraction or combine GPU agents with
CPU matching.

\textbf{PyMarketSim}~\cite{mascioli2024} (ACM ICAIF 2024) is a Python discrete-time simulator whose
tick (observe, act, clear, advance) matches ours. Its order book takes 1.9--4.7\,$\mu$s per
operation from $10^3$ to $10^6$ orders, about 300{,}000 book operations or 15{,}000 agent actions
per second, but it is serial, simulates one asset and uses 25 agents in its experiments.

\subsection{GPU and high-performance ABM frameworks}
\textbf{FLAME GPU 2}~\cite{richmond2023} (Software: Practice and Experience 2023) compiles agent
functions to CUDA at run time, stores agents as structures of arrays, and passes messages through
spatial, array or bucket structures. It runs up to 200 million agents on a V100 and 1M agents in
about 3\,ms per step. Its performance depends on spatial or keyed locality; agents competing for
one resource need repeated bidding sub-models, atomics are non-deterministic, and one model is
limited to one GPU.

\textbf{Agents.jl}~\cite{datseris2024} (SIMULATION 2024, online 2022) is a Julia ABM library that
runs single models about 2--125$\times$ faster than MASON, NetLogo and Mesa in the authors'
benchmarks. It does not parallelise within a model, because agents mutate shared state, and has no
GPU support: the same shared-state conflict that an order book creates.

\subsection{Matching engines and hardware-accelerated trading}
\textbf{CoinTossX}~\cite{jericevich2022} (SoftwareX 2022) is an open-source Java matching engine
modelled on the Johannesburg Stock Exchange, with price--time priority and messaging over UDP. It
reaches 106--248\,ns 90th-percentile matching latency on a 16-core server, but throughput falls from
2{,}287 to 621--989 orders per second as client--stock pairs are added, and it is engineered for
latency rather than for parallel scaling.

\textbf{LightTrader}~\cite{yoo2023} (IEEE HPCA 2023) is a heterogeneous high-frequency trading system:
an FPGA parses market data and keeps the top of a local order book, while custom deep-learning
accelerators run price-prediction models, with a scheduler that picks batch size and voltage and
frequency under a power budget. It is 13.9$\times$ faster than a V100-based system for AI
processing, but scaling saturates as accelerators share the power budget. It accelerates the
\emph{decision} stage for one trader, the analogue of our GPU phase.

\subsection{Distributed ABM and reproducible parallel random numbers}
\textbf{Repast4Py}~\cite{collier2022} (Winter Simulation Conference 2022) distributes Python ABMs
with MPI, partitioning space across ranks with ghost regions. A 3.6M-agent random walk speeds up
4.39$\times$ from 36 to 288 ranks (efficiency 0.94 to 0.55), while a network model reaches only
1.75$\times$ (efficiency 0.22): communication dominates at higher rank counts. Whether results are
identical across rank counts is not addressed.

\textbf{SIMCoV-GPU}~\cite{leyba2024} (ACM HPDC 2024) ports a spatial infection ABM to multiple GPUs
with CUDA and UPC++, reaching 11.9$\times$ over 512 CPU cores with 16 GPUs. The GPU wins only on
large enough problems (at 64 GPUs on a small grid it is 0.85$\times$), and because CPU and GPU
random number generators differ, results agree only statistically (within about 1\%), not bit for
bit.

\textbf{OpenRAND}~\cite{khan2024} (SoftwareX 2024) provides counter-based generators (Philox,
Threefry and others) keyed by a logical entity rather than a thread, so streams give the same
output on any thread count or device; it is 1.8$\times$ faster than cuRAND in a particle
benchmark. It supplies reproducible random numbers but not reproducibility of a whole simulation,
which also needs deterministic arithmetic and event ordering.

\subsection{Summary, research gaps and motivation}
Table~\ref{tab:lit} summarises the reviewed work. Four gaps motivate this project:
\begin{enumerate}
\item \textbf{Market simulators are built for fidelity, not parallel characterisation.}
  ABIDES-Gym, PyMarketSim and CoinTossX run essentially serially; JAX-LOB parallelises only across
  independent books. None splits the tick into a parallel decision phase and a serial matching
  phase and relates the measured speedup to Amdahl's law.
\item \textbf{The serial matching core is acknowledged but not measured.} No study reports what
  share of tick time matching takes as agents grow, or the gap between kernel and end-to-end
  speedup.
\item \textbf{GPU and distributed ABM frameworks rely on locality.} FLAME GPU 2, Repast4Py and
  SIMCoV-GPU exploit spatial neighbourhoods and halos; Agents.jl avoids in-model parallelism. An
  order book, where every agent writes into one strictly ordered structure per instrument, fits
  none of them.
\item \textbf{Bit-identical heterogeneous results are rare.} SIMCoV-GPU accepts statistical
  agreement and FLAME GPU 2 has non-deterministic atomics; OpenRAND makes only the random numbers
  reproducible. No reviewed work shows a bit-identical event log across CUDA, OpenMP and MPI.
\end{enumerate}
We therefore build a simulator that partitions by instrument, keeps the GPU for decisions only,
produces identical trades on every back end, and measures where its speedup stops and why.

\begin{widetable}
\centering
\caption{Comparison of recent related work (2021--2024).}
\label{tab:lit}
\footnotesize
\begin{tabularx}{\textwidth}{@{}p{2.3cm} X X X p{2.0cm}@{}}
\toprule
Work & Approach / platform & Advantages & Limitations & Key result\\
\midrule
ABIDES-Gym~\cite{amrouni2021} & discrete-event kernel, Python, serial & high fidelity, latency model, RL interface & single-threaded; no run times reported & --\\
JAX-LOB~\cite{frey2023} & JAX on GPU, \texttt{vmap} over books & thousands of books in parallel; RL on GPU & serial within a book; fixed capacity & 2.6\,$\mu$s/msg per book; 7$\times$ RL\\
PyMarketSim~\cite{mascioli2024} & Python discrete-time tick & simple, RL and game-theory tools & serial; one asset; 25 agents & 300k book ops/s\\
FLAME GPU 2~\cite{richmond2023} & CUDA, SoA, spatial messages & 200M agents on one GPU & needs locality; non-det.\ atomics; 1 GPU per model & 1M agents in 3\,ms/step\\
Agents.jl~\cite{datseris2024} & Julia, CPU & concise and fast single models & no in-model parallelism; no GPU & 2--125$\times$ vs other ABM tools\\
CoinTossX~\cite{jericevich2022} & Java engine, UDP & sub-$\mu$s matching latency & throughput falls with instruments & 106--248\,ns latency\\
LightTrader~\cite{yoo2023} & FPGA + DL accelerators & fast decision stage, energy-aware & power-limited scaling; one trader & 13.9$\times$ vs V100 system\\
Repast4Py~\cite{collier2022} & Python + MPI, spatial partitions & HPC scale from Python & communication-bound at high rank counts & 4.39$\times$, 36$\to$288 ranks\\
SIMCoV-GPU~\cite{leyba2024} & CUDA + UPC++, multi-GPU & 11.9$\times$ over 512 cores & GPU wins only at large sizes; not bitwise reproducible & 11.9$\times$ (16 GPUs)\\
OpenRAND~\cite{khan2024} & counter-based RNG library & reproducible on any back end & RNG only, not whole simulation & 1.8$\times$ vs cuRAND\\
\bottomrule
\end{tabularx}
\end{widetable}
